\documentclass[10pt,a4paper]{amsart}
\usepackage[margin=19mm]{geometry}
\usepackage{amsmath,amssymb,mathtools}
\usepackage[scr=rsfs,cal=euler]{mathalfa}
\usepackage{xfrac}
\usepackage[unicode,hidelinks,bookmarks=false]{hyperref}
\newcommand{\ie}{i.e.\ }
\newcommand{\cf}{cf.\ }
\newcommand{\resp}{respectively\ }
\newcommand{\Subsection}{\subsection}
\providecommand{\nobreakdash}{\nobreak\hbox{-}}
\setlength{\emergencystretch}{2em}
\multlinegap0pt
\usepackage{amssymb}
\usepackage{tikz}
\usepackage{bm}
\usepackage{algorithm}\usepackage{algpseudocode}
\usepackage{float}
\usepackage[shortlabels]{enumitem}
\newtheorem{theorem}{Theorem}
\newcommand{\Ltwo}{{L^2}}
\newcommand{\N}{\field{N}}
\newcommand{\Sone}{{S^1}}
\newcommand{\Z}{\field{Z}}
\newcommand{\bfF}{{\bs F}}
\newcommand{\bfa}{{\bs a}}
\newcommand{\bfc}{{\bs c}} 
\newcommand{\bfe}{{\bs e}}
\newcommand{\bfy}{{\bs y}} 
\newcommand{\bfzero}{{\bs 0}}
\newcommand{\bs}{\boldsymbol} 
\DeclareMathOperator{\dif}{d\!}  
\newcommand{\dy}{\, \dif \bfy} 
\newcommand{\field}[1]{{\mathbb{#1}}}
\newcommand{\ph}{\,\cdot\,}
\newcommand{\rme}{\mathrm{e}}
\newcommand{\rmi}{\mathrm{i}} 
\newcommand{\yhat}{\widehat{\bfy}}
\title{Direct reconstruction for acoustic inverse Born scattering}
\author{Nuutti Hyv\"onen, Lisa Sch\"atzle}
\date{Source: arXiv:2606.05977v1 (CC BY 4.0)}
\begin{document}
\maketitle

\noindent\emph{Source and licence.} Direct reconstruction for acoustic inverse Born scattering. Version arXiv:2606.05977v1 (CC BY 4.0), distributed by arXiv under the Creative Commons Attribution 4.0 International licence, http://creativecommons.org/licenses/by/4.0/, as recorded in the arXiv licence metadata for this exact version and shown on the abstract page https://arxiv.org/abs/2606.05977v1. Full licence notice: \texttt{licenses/arxiv-2606.05977-CC-BY-4.0.txt}.\par\smallskip

\noindent\textit{Editorial scope.} Counted unit: the article's theorem Theorem:main (source0.tex lines 1170--1180). The article's complete original proof of it is delivered with it (source0.tex lines 1182--1189, one \texttt{proof} environment of 8 lines). No mathematics is written, repaired or re-derived: the statement and the proof are the article's own lines, in the article's own notation, and no retained line is substituted or re-flowed.  Retained uncounted as the source-local context the proof needs: lines 810--816; lines 826--830; lines 836--840; lines 976--981; lines 1155--1158.  Nothing was omitted from any retained span, so the package carries no omission record.  No retained line contains a citation, so the wrapper carries no bibliography and no bibliography entry.  No retained line contains a figure environment, an image-inclusion command or drawing code, so no image is delivered and the package declares no asset dependency.  Editorial formatting only, in addition: an AMS \texttt{amsart} preamble that restates the article's own declarations and macros used by the retained lines, namely the environments \texttt{theorem} on separate counters, together with the standard packages those lines need.  The retained environments print in this document as Theorem 1 in order of appearance, whereas the article prints the same items with the numbering of its own full document.  The complete native source of the article as distributed in the arXiv e-print for this exact version is bundled unchanged as \texttt{sources/202150/source0.tex}.
\begin{align}
	a_{m,n} &\,:=\, \langle F_B\bfe^\bfc_n,\bfe^\bfc_m\rangle_{L^2(\Sone)} \notag\\
    &\,=\, 2\pi \kappa^2\rmi^{n-m} \int_{B_R(\bfzero)} q(\bfy+\bfc)\rme^{\rmi(n-m)\arg\bfy}J_m(\kappa|\bfy|)J_n(\kappa|\bfy|)\dy \notag\\
	&\,=\, (2\pi)^{3/2} \kappa^2\rmi^{n-m} \int_{B_R(\bfzero)} q(\bfy+\bfc)\overline{\bfe_{m-n}(\yhat)}J_m(\kappa|\bfy|)J_n(\kappa|\bfy|)\dy \notag\\
    &\,=\, (2\pi)^{3/2} (\kappa R)^2\rmi^{n-m} \int_{B_1(\bfzero)} q(R\bfy+\bfc)\overline{\bfe_{m-n}(\yhat)}J_m(\kappa R|\bfy|)J_n(\kappa R|\bfy|)\dy \,,
    \label{eq:exp_coeff_FF}
\end{align}

\begin{equation}
\label{eq:def_Psi}
    \Psi_{j,k}(\bfy):=\bfe_j(\yhat)R_k^{|j|}(|\bfy|),
    \qquad  \bfy\in B_1(\bfzero)\,.
\end{equation}

\begin{equation}
\label{eq:exp_qc}
	q(R(\ph)+\bfc) \,=\, \sum_{j\in\Z} q_j \,=\, \sum_{j\in\Z}\sum_{k=0}^\infty c_{j,k} \Psi_{j,k}
	\qquad \text{for } c_{j,k}\,:=\, \big \langle q(R(\ph)+\bfc), \Psi_{j,k} \big \rangle_{\Ltwo(B_1(\bfzero))}\,,
\end{equation}

\begin{equation}
\label{eq:def_radial_basis}
  R_k^{|j|} \,:=\, \frac{\widetilde R_k^{|j|}}{\big \|\widetilde R_k^{|j|} \big \|_{L^2_r(0,1)}} 
    \qquad \text{with} \  \widetilde R_k^{|j|}
    \,:=\, P_{k+\lceil|j|/2\rceil}^{|j|} - \sum_{m=0}^{k-1} \big \langle P_{k+\lceil|j|/2\rceil}^{|j|},R_{m}^{|j|} \big \rangle_{L^2_r(0,1)}R_{m}^{|j|}
\end{equation}

\begin{equation}
\label{eq:infinite_system}
	\bfF^{j} \bfc^j\,=\, \bfa^j
\end{equation}

\begin{theorem}
\label{Theorem:main}
Denote by $(\Psi_{j,k})_{j\in\Z,k\in\N_0}$ the orthonormal system for $\Ltwo(B_1(\bfzero))$ defined by \eqref{eq:def_Psi} and \eqref{eq:def_radial_basis}, and let $(a_{m,n})_{m,n\in\Z}$ be the given expansion coefficients of the Born far field operator as in \eqref{eq:exp_coeff_FF}.
Then, the expansion coefficients of the shifted and scaled contrast $q(R(\ph)+\bfc)$ with respect to $(\Psi_{j,k})_{j\in\Z,k\in\N_0}$, as given in \eqref{eq:exp_qc}, can be computed separately for each $j \in \Z$ via a recursion with respect to $k$:
\begin{equation}
\label{eq:forward_substitution}
	c_{j,k} \,=\, \frac 1 {(2\pi)^{3/2}(\kappa R)^2(-\rmi)^j \big\| \widetilde{R}_{k}^{|j|} \big \|_{L^2_r(0,1)}} a_{k+\lceil j/2\rceil, k-\lfloor j/2\rfloor} 
	- \sum_{i=0}^{k-1} \frac{\big \langle R_{i}^{|j|}, P_{k+\lceil |j|/2\rceil}^{|j|} \big \rangle_{L^2_r(0,1)}}{\big\| \widetilde{R}_{k}^{|j|} \big\|_{L^2_r(0,1)} } \, c_{j, i}
\end{equation}
for $k \in \N_0$. Here, $(\widetilde{R}_{k}^{|j|})_{k \in \N_0}$ are the unnormalized orthogonal basis functions from \eqref{eq:def_radial_basis} and the sum in \eqref{eq:forward_substitution} is viewed to be empty for $k=0$.
\end{theorem}

\begin{proof}
Fix $j \in \Z$. Solving the lower triangular system \eqref{eq:infinite_system} via forward substitution directly gives
\begin{equation*}
	c_{j,k} \,=\, \frac 1 {(2\pi)^{3/2}(\kappa R)^2(-\rmi)^j \big \langle R_{k}^{|j|}, P_{k+\lceil |j|/2\rceil}^{|j|} \big \rangle_{L^2_r(0,1)}} a_{k+\lceil j/2\rceil, k-\lfloor j/2\rfloor} 
	- \sum_{i=0}^{k-1} \frac{\big \langle R_{i}^{|j|}, P_{k+\lceil |j|/2\rceil}^{|j|} \big \rangle_{L^2_r(0,1)}}{\big \langle R_{k}^{|j|}, P_{k+\lceil |j|/2\rceil}^{|j|}\big \rangle_{L^2_r(0,1)}} c_{j, i} \, 
\end{equation*}
for $k \in \N_0$. Representing $P_{k+\lceil |j|/2\rceil}^{|j|}$ as in \eqref{eq:def_radial_basis} and utilizing the orthogonality of $(R_{k}^{|j|})_{k \in \N_0}$ allows to replace the inner product in the denominator by the norm of $\widetilde{R}_{k}^{|j|}$, which is nonzero by construction. This completes the proof. 
\end{proof}

\end{document}
